\documentclass[10pt,a4paper]{amsart}
\usepackage[margin=19mm]{geometry}
\usepackage{amsmath,amssymb,mathtools}
\usepackage[scr=rsfs,cal=euler]{mathalfa}
\usepackage{xfrac}
\usepackage[unicode,hidelinks,bookmarks=false]{hyperref}
\newcommand{\ie}{i.e.\ }
\newcommand{\cf}{cf.\ }
\newcommand{\resp}{respectively\ }
\newcommand{\Subsection}{\subsection}
\providecommand{\nobreakdash}{\nobreak\hbox{-}}
\setlength{\emergencystretch}{2em}
\multlinegap0pt
\usepackage{bm}
\usepackage{bbm}
\usepackage{dsfont}
\usepackage{xspace}
\usepackage{cancel}
\usepackage{mathtools}
\usepackage{amsthm}
\makeatletter
\@ifundefined{lemma}{\newtheorem{lemma}{Lemma}}{}
\@ifundefined{proposition}{\newtheorem{proposition}{Proposition}}{}
\makeatother
\makeatletter
\def\E{\mathbb E}
\def\R{\mathbb R}
\expandafter\ifx\csname assumptionp\endcsname\relax
\newenvironment{assumptionp}[1]{
\renewcommand\theassumptionalt{#1}
\assumptionalt
}{\endassumptionalt}
\fi
\expandafter\ifx\csname conditionp\endcsname\relax
\newenvironment{conditionp}[1]{
\renewcommand\theconditionalt{#1}
\conditionalt
}{\endconditionalt}
\fi
\makeatother
\title[Backward stochastic differential equations with nonlinear Yo]{Backward stochastic differential equations with nonlinear Young drivers I}
\author{Jian Song, Huilin Zhang, Kuan Zhang}
\date{Source: arXiv:2504.18632v3 (CC BY 4.0)}
\begin{document}
\maketitle

\noindent\emph{Source and licence.} Backward stochastic differential equations with nonlinear Young drivers I. Version arXiv:2504.18632v3 (CC BY 4.0), distributed under the Creative Commons Attribution 4.0 International licence, as recorded in the arXiv licence metadata for this exact version and shown on the abstract page https://arxiv.org/abs/2504.18632v3. Full rights notice: licenses/arxiv-2504.18632-CC-BY-4.0.txt.\par\smallskip

\noindent\textit{Editorial scope.} Counted unit: the Proposition whose source label is \texttt{prop:Ito's formula} in the article (source0.tex lines 435--442, source label \texttt{prop:Ito's formula}), delivered together with the article's own complete original proof of it (source0.tex lines 444--612, one \texttt{proof} environment of 169 lines). No mathematics is written, repaired or re-derived: statement and proof are the article's own text line for line. Required context retained uncounted: prop:bounds (lines 350--368); e:p-var1 (lines 362--366); Y-decomp (lines 429--431); lem:BDG (lines 2135--2142); lem:sewing (lines 2304--2317); lem:control's property (lines 2320--2324); e:inequality-estimate-Young (lines 2369--2373). Every definition, display and label of those lines is retained unchanged. Omitted from this package and not redistributed: 1--349, 367--428, 432--434, 443--443, 613--2134, 2143--2303, 2318--2319, 2325--2368, 2374--2451. Nothing inside a retained excerpt is omitted or altered: every retained body is delivered exactly as the article's lines are written, so no omission receipt of any kind accompanies this package. Citations: the retained excerpts carry no citation command at all, so no bibliography entry of the article is transcribed and no reference is redistributed. Reference adaptation: none was needed and none was made; every reference inside the delivered unit resolves inside it, so no unresolved reference or citation is produced. Printed slips kept literal: no printed word, symbol, hypothesis or number of the counted statement or of its proof was altered. Layout and class: the article's own class is \texttt{amsart} with packages including amsthm, amssymb, amsmath, epsfig, graphics, appendix, bm, mathrsfs, bbm, float, graphicx, hyperref, xcolor, stmaryrd; the wrapper is the lane's pinned \texttt{amsart} editorial wrapper at 10pt on A4, which restates the theorem-like environments the retained text needs and adds the article's own macro definitions that the retained text needs (\texttt{\textbackslash E}, \texttt{\textbackslash R}), with extra wrapper packages (bm, bbm, dsfont, xspace, cancel, mathtools). The delivered numbering is the unit's own. Assets: no retained span contains a figure environment, a graphics-inclusion command, a PDF-page inclusion, a table or a TikZ picture; no image file is copied into this package, the package declares no asset dependency and no image file is redistributed with it. Licence: version arXiv:2504.18632v3 of this article is distributed under the Creative Commons Attribution 4.0 International licence, as recorded in the arXiv licence metadata for this exact version and shown on the abstract page https://arxiv.org/abs/2504.18632v3. A full rights notice is delivered at licenses/arxiv-2504.18632-CC-BY-4.0.txt. Characters: every retained excerpt is pure ASCII, so the delivered engine, XeLaTeX with its default Latin Modern text font, reports no missing character.
\begin{proposition}\label{prop:bounds}
Let $ \tau\in(0, 1]$, $ \lambda\in(0,1]$. 
Assume $\eta\in C^{\tau,\lambda}_{\mathrm{loc}}([0,T]\times \mathbb{R}^{d})$, $x\in C^{p_{1}\text{-}\mathrm{var}}([0,T];\mathbb{R}^{d})$, and $y\in C^{p_{2}\text{-}\mathrm{var}}([0,T];\mathbb{R})$ with $\tau +\frac{\lambda}{p_{1}}>1$  and $\tau + \frac{1}{p_{2}}>1$.  For $0\le a\le b\le T$, the following integral is well-defined,		
\[\int_{a}^{b} y_r \eta(dr,x_r) := \mathcal{A}_{a,b} := \mathcal A(b)-\mathcal A(a),\] 
where $\mathcal{A}:[0,T]\to \R$ is the unique function associated with $A(s,t)= y_s\big(\eta(t,x_s) - \eta(s,x_s)\big)$ by Lemma~\ref{lem:sewing}. Moreover, if we assume $\eta\in C^{\tau,\lambda;\beta}([0,T]\times\mathbb{R}^{d})$ for some $\beta\ge 0$, 
\begin{equation}\label{e:p-var2}
\begin{aligned}
\left\|\int_{\cdot}^{t}y_{r}\eta(dr,x_{r})\right\|_{\frac{1}{\tau}\text{-}\mathrm{var};[s,t]}&\le \frac{2^{\delta}}{1-2^{-\delta}}\|\eta\|_{\tau,\lambda;\beta} |t-s|^{\tau} \bigg(\Big(1+\|x\|^{\beta}_{\infty;[s,t]}\Big)\|x\|^{\lambda}_{p_{1}\text{-}\mathrm{var};[s,t]}\|y\|_{\infty;[s,t]}\\ 
&\quad +  \Big(1+\|x\|^{\beta+\lambda}_{\infty;[s,t]}\Big)\Big(\|y\|_{\infty;[s,t]}+\|y\|_{p_{2}\text{-}\mathrm{var};[s,t]}\Big)\bigg),
\end{aligned}
\end{equation}
and if we assume $\eta \in C^{\tau,\lambda}([0,T]\times\mathbb{R}^{d})$, 
\begin{equation}\label{e:p-var1}
\begin{aligned} 
\left\|\int_{\cdot}^{t}y_{r}\eta(dr,x_{r})\right\|_{\frac{1}{\tau}\text{-}\mathrm{var};[s,t]}\le \frac{2^{\delta}}{1-2^{-\delta}}\|\eta\|_{\tau,\lambda} |t-s|^{\tau} \bigg(\Big(1+\|x\|^{\lambda}_{p_{1}\text{-}\mathrm{var};[s,t]}\Big)\|y\|_{\infty;[s,t]} + \|y\|_{p_{2}\text{-}\mathrm{var};[s,t]} \bigg),
\end{aligned} 
\end{equation}
for all $0\le s\le t\le T$, where $\delta = \min\left\{\tau+\frac{1}{p_{2}}-1,\tau+\frac{\lambda}{p_{1}}-1\right\}$.
\end{proposition}

\begin{equation}\label{Y-decomp}	
Y^{i}_{t} = Y^{i}_{0} - \int_{0}^{t} f^{i}_{r} dr - \sum_{j=1}^{M}\int_{0}^{t}g^{i}_{j}(r) \eta_{j}(dr,X_{r}) + \int_{0}^{t}Z^{i}_{r}dW_{r} ,\ i=1,...,n,\ t\in[0,T],
\end{equation}

\begin{lemma}\label{lem:BDG}
Assume $p>2$, $k\ge 1$, and $t\in[0,T]$. Let $M$ be a continuous local martingale such that $\E[\langle M\rangle_{t,T}^{k/2}]<\infty$.  Then, there exists a positive number $C$ depending only on $p$ and $k$ such that
\[C^{-1}\E_{t}[\langle M\rangle_{t,T}^{\frac k2}]\le \mathbb{E}_{t}[\|M\|^{k}_{p\text{-}\mathrm{var};[t,T]}] \le C\E_{t}[\langle M\rangle_{t,T}^{\frac k2}].\]
If we further assume that $\mathbb{E}[\|M\|^{k}_{\infty;[t,T]}]<\infty$ (hence $M$ is a martingale), the above result, combined with the lower bound of the classical BDG inequality, yields 
\[\mathbb{E}_{t}[\|M\|^{k}_{p\text{-}\mathrm{var};[t,T]}] \lesssim_{p,k} \mathbb{E}_{t}[\|M\|^{k}_{\infty;[t,T]}].\]
In addition, Doob's inequality yields, for $k>1$,
\[\mathbb{E}_{t}[\|M\|^{k}_{p\text{-}\mathrm{var};[t,T]}] \lesssim_{p,k} \mathbb{E}_{t}[|M_T|^{k}].\]
\end{lemma}

\begin{lemma}\label{lem:sewing} 
Let $A: \{(s, t); 0 \leq s \le t \leq T\} \rightarrow \mathbb{R}$ be a continuous function satisfying
\begin{equation*}
|\delta A_{s,u,t}| \le \sum_{i = 1}^{l} w_{i}(s,t)^{1+\varepsilon_{i}},\quad 0\le s\le u\le t \le T,
\end{equation*}
for some $\varepsilon_{i}>0,\ i=1,2,...,l$, where $w_{1}$, $w_{2}$,..., $w_{l}$ are controls and $$\delta A_{s,u,t} := A(s,t) - A(s,u) - A(u,t).$$ 

Then, there exists a unique function $\mathcal{A}: [0,T]\rightarrow \mathbb{R}$ satisfying that
\begin{equation*}
|\mathcal{A}_{s,t} - A(s,t)|\le \frac{l^{\varepsilon_{0}}}{1-2^{-\varepsilon_{0}}}\sum_{i = 1}^{l} w_{i}(s,t)^{1+\varepsilon_{i}},\quad  0 \le s \le t \le T,
\end{equation*}
where $\mathcal{A}_{s,t} := \mathcal{A}(t) - \mathcal{A}(s)$, $\varepsilon_{0} := \min\{\varepsilon_{i}, i=1, \dots, l\}$,
and \[\mathcal{A}_{s,t}=\lim_{|\pi|\downarrow 0} \sum\limits_{[t_{i},t_{i+1}]\in\pi} A(t_{i}, t_{i+1}).\] Here, $|\pi|:=\max_{i}|t_i-t_{i-1}|$  is the maximum length of subintervals of the partition $\pi:=\{[t_i,t_{i+1}] ;\  s=t_0<t_1<\cdots<t_n=t\}$, which is also called the mesh of $\pi$. 
\end{lemma}

\begin{lemma}\label{lem:control's property}
Suppose that for every $0\le s \le t \le T$, $|\Gamma(s,t)|\le w(s,t)^{\theta}$ for some $\theta > 1$, where $w$ is a control function.  Then, we have
\[z_{r} := \lim_{|\pi|\downarrow 0}\sum_{[t_{i},t_{i+1}]\in\pi} \Gamma(t_{i}\land r,t_{i+1}\land r)=0,\ \ \ \text{for all }r\in[0,T].
\] 
\end{lemma}

\begin{equation}\label{e:inequality-estimate-Young}
\begin{aligned}
\left|\int_{s}^{t} y_{r} \eta(dr,x_{r}) - A(s,t)\right| &\le \frac{2^{\delta}}{1-2^{-\delta}}\Big(\hat{w}_1(s,t)^{1+\delta}+\hat{w}_2(s,t)^{1+\delta}\Big).
\end{aligned}
\end{equation}

\begin{proposition}[It\^o's formula]\label{prop:Ito's formula} Let $ \tau\in(\frac12, 1]$, $ \lambda\in(0,1]$, $p_1\ge 1$, and  $p_2\ge1$ be parameters such that $\tau + \frac{\lambda}{p_{1}}>1$ and $\tau + \frac{1}{p_{2}}>1$.  Suppose $Y$ satisfies \eqref{Y-decomp}. Then, for any $\phi(t,y)\in C^{\gamma,1}_{\mathrm{loc}}([0,T]\times\mathbb{R}^{n})$ with $\gamma\in (\frac12, 1],$ having continuous spatial partial derivatives $\nabla_{y}\phi$ and $\nabla^{2}_{yy}\phi$, we have
\begin{equation*}
\begin{aligned}
\phi(t,Y_{t}) &=   \phi(0,Y_{0}) + \int_{0}^{t}\ \phi(dr,Y_{r}) -\int_{0}^{t}\nabla_{y}\phi(r,Y_{r})^{\top} f_{r} dr - \sum_{i=1}^{n}\sum_{j=1}^{M}\int_{0}^{t} \partial_{y^i}\phi(r,Y_{r}) g^{i}_{j}(r)\eta_{j}(dr,X_{r})\\
&\quad + \int_{0}^{t}  \nabla_{y}\phi(r,Y_{r})^{\top} Z_{r} dW_{r} + \frac{1}{2}\int_{0}^{t}\text{tr}\left\{\nabla^{2}_{yy}  \phi(r,Y_{r})Z_{r}Z^{T}_{r}\right\} dr,\ t\in[0,T].
\end{aligned}
\end{equation*}
\end{proposition}

\begin{proof}[Proof of Proposition~\ref{prop:Ito's formula}]
\makeatletter\def\@currentlabelname{the proof}\makeatother
\label{proof:Ito}
We only prove the case when $n=d=M=1$; the general case can be proved similarly. Since $\gamma,\tau>\frac{1}{2}$, there exists a real number $p_{3} > 2$  such that $\gamma + \frac{1}{p_{3}} > 1$ and $\tau + \frac{1}{p_{3}}>1$. Furthermore, Since $p_3>2,$ by Proposition~\ref{prop:bounds} and Lemma~\ref{lem:BDG}, we have $Y\in C^{p_{3}\text{-var}}([0,T];\mathbb{R})$ a.s. In addition, since $\phi\in C^{\gamma,1}_{\text{loc}}([0,T]\times\mathbb{R})$ and $\gamma + \frac{1}{p_{3}}>1$, we have, for almost all paths of $Y$,  $\int_{0}^{t}\phi(dr,Y_{r})$ is well defined in the sense of Proposition~\ref{prop:bounds}. Also, by our assumption, $\nabla_{y}\phi(\cdot,x)$ is $\gamma$-H\"older continuous in $[0,T]$ for every fixed $ x\in \mathbb{R}$. 

Let $\left\{\tau_{K}\right\}_{K=1}^{\infty}$   be a sequence of stopping times such that $\tau_{K} \uparrow T$ a.s., and the following terms are all less than $K$ (here $X^{\tau_{K}}_{\cdot} := X_{\cdot\land\tau_{K}}$; $g^{\tau_{K}}$ and $Y^{\tau_{K}}$ are defined in the same manner):
\begin{equation*}
\begin{aligned}
&\|X^{\tau_{K}}_{\cdot}\|_{\infty;[0,T]},\ \|X^{\tau_{K}}_{\cdot}\|_{p_{1}\text{-var};[0,T]},\ \|g^{\tau_{K}}_{\cdot}\|_{\infty;[0,T]},\ \|g^{\tau_{K}}_{\cdot}\|_{p_{2}\text{-var};[0,T]},\ \|Y^{\tau_{K}}_{\cdot}\|_{\infty;[0,T]},\ \|Y^{\tau_{K}}_{\cdot}\|_{p_{3}\text{-var};[0,T]},\\
&\int_{0}^{\tau_{K}}|Z_{r}|^{2}dr,\ \Big\|\int_{0}^{\cdot\land\tau_{K}}Z_{r}dW_{r}\Big\|_{p_{3}\text{-var};[0,T]},\ \int_{0}^{\tau_{K}}|f_{r}|dr,\ \text{and}\ \Big\|\int_{0}^{\cdot\land \tau_{K}}g_{r}\eta(dr,X_{r})\Big\|_{\frac{1}{\tau}\text{-var};[0,T]}.
\end{aligned}
\end{equation*}
Denote $B^{K}:=\{x\in\R;|x|\le K\}$ and  $\|\eta\|^{(K)}_{\tau,\lambda} := \|\eta\|^{(B^{K})}_{\tau,\lambda}$ for simplicity; the notation $\|\phi\|^{(K)}_{\gamma,1}$ is defined in the same manner. Hence, letting $\pi$ be a partition on $[0,t]$, we have 
\begin{equation*}
\begin{aligned}
\Bigg(&\sum_{[t_{i},t_{i+1}]\in\pi}\left|\partial_{y}\phi(t_{i+1},Y^{\tau_{K}}_{t_{i+1}}) - \partial_{y}\phi(t_{i},Y^{\tau_{K}}_{t_{i}})\right|^{p_{3}}\Bigg)^{\frac{1}{p_{3}}}\\
&\le \left(\sum_{i}\left|\partial_{y}\phi(t_{i+1},Y^{\tau_{K}}_{t_{i+1}}) - \partial_{y}\phi(t_{i+1},Y^{\tau_{K}}_{t_{i}})\right|^{p_{3}}\right)^{\frac{1}{p_{3}}} + \left(\sum_{i}\left|\partial_{y}\phi(t_{i+1},Y^{\tau_{K}}_{t_{i}}) - \partial_{y}\phi(t_{i},Y^{\tau_{K}}_{t_{i}})\right|^{p_{3}}\right)^{\frac{1}{p_{3}}}\\
&\le \sup_{r\in[0,T],|y|\le K} |\partial^{2}_{yy}\phi(r,y)|\left(\sum_{i}\left|Y^{\tau_{K}}_{t_{i+1}} - Y^{\tau_{K}}_{t_{i}}\right|^{p_{3}}\right)^{\frac{1}{p_{3}}} + \|\phi\|^{(K)}_{\gamma,1} \left(\sum_{i}\left|t_{i+1}-t_{i}\right|^{\gamma \cdot p_{3}}\right)^{\frac{1}{p_{3}}}.
\end{aligned}
\end{equation*}
Therefore, noting that $\gamma\cdot p_{3} > p_{3} - 1 > 1,$  $\partial_{y}\phi(\cdot,Y^{\tau_{K}}_{\cdot})$ has finite $p_{3}$-variation. Hence, since $\tau + \frac{1}{p_{3}}>1$, the nonlinear Young integral $\int_{a}^{b}\partial_{y}\phi(r,Y^{\tau_{K}}_{r})\eta(dr,X_{r})$ is well-defined for $0\le a <b \le T$. Note that
\begin{equation}\label{decomI}	
\begin{aligned}
\phi(t,Y^{\tau_{K}}_{t}) - \phi(0,Y_{0}) &= \sum_{[t_{i},t_{i+1}]\in \pi}\left[\phi(t_{i+1},Y^{\tau_{K}}_{t_{i+1}}) - \phi(t_{i},Y^{\tau_{K}}_{t_{i}})\right]  =: I_{1} + I_{2},
\end{aligned} 
\end{equation}
where \[I_{1} := \sum_{[t_{i},t_{i+1}]\in\pi}\left[\phi(t_{i+1},Y^{\tau_{K}}_{t_{i+1}}) - \phi(t_{i+1},Y^{\tau_{K}}_{t_{i}})\right]  \text{ and } I_{2} := \sum_{[t_{i},t_{i+1}]\in\pi}\left[\phi(t_{i+1},Y^{\tau_{K}}_{t_{i}}) - \phi(t_{i},Y^{\tau_{K}}_{t_{i}})\right]
.\]

We first deal with the term $I_{2}$, which can be written as	follows,	
\begin{equation*}
\begin{aligned}
I_{2} & = \sum_{[t_{i},t_{i+1}]\in\pi} \left[\int_{t_{i}}^{t_{i+1}}\phi(dr,Y^{\tau_{K}}_{r}) + \phi(t_{i+1},Y^{\tau_{K}}_{t_{i}}) - \phi(t_{i},Y^{\tau_{K}}_{t_{i}}) - \int_{t_{i}}^{t_{i+1}}\phi(dr,Y^{\tau_{K}}_{r})\right]\\
& = \int_{0}^{t}\phi(dr,Y^{\tau_{K}}_{r}) +\sum_{[t_{i},t_{i+1}]\in\pi}\left[ \phi(t_{i+1},Y^{\tau_{K}}_{t_{i}}) - \phi(t_{i},Y^{\tau_{K}}_{t_{i}}) - \int_{t_{i}}^{t_{i+1}}\phi(dr,Y^{\tau_{K}}_{r})\right].
\end{aligned}
\end{equation*}
By \eqref{e:inequality-estimate-Young} we have 	
\begin{align*}	
\left|\phi(t_{i+1},Y^{\tau_{K}}_{t_{i}}) - \phi(t_{i},Y^{\tau_{K}}_{t_{i}}) - \int_{t_{i}}^{t_{i+1}}\phi(dr,Y^{\tau_{K}}_{r})\right| \lesssim_{\gamma,p_{3}} \|\phi\|^{(K)}_{\gamma,1}(t_{i+1}-t_{i})^{\gamma} \|Y^{\tau_{K}}_{\cdot}\|_{p_{3}\text{-var};[t_{i},t_{i+1}]}.
\end{align*}
Noting that $\gamma + \frac{1}{p_{3}}>1$, by Lemma~\ref{lem:control's property}, it follows that
\begin{equation}\label{e:I2}
\lim_{|\pi|\downarrow 0} I_{2} = \int_{0}^{t}\phi(dr,Y^{\tau_{K}}_{r})\ \ \ \ a.s.	
\end{equation}

In the remaining part of the proof, we analyze $I_{1}$, which is decomposed as follows,
\begin{equation}\label{e:I1}	\begin{aligned}
I_{1} &= \sum_{[t_{i},t_{i+1}]\in \pi}\left[\int_{t_{i}}^{t_{i+1}}\partial_{y}\phi(r,Y^{\tau_{K}}_{r})dY^{\tau_{K}}_{r} + \frac{1}{2}\int_{t_{i}}^{t_{i+1}}\partial^{2}_{yy}\phi(r,Y^{\tau_{K}}_{r})|Z_{r}|^{2}\mathbf{1}_{[0,\tau_{K}]}(r) dr\right]\\
& \qquad + J_{1} + J_{2} + J_{3},
\end{aligned}
\end{equation}
where $dY^{\tau_{K}}_{r} := -f_{r}\mathbf{1}_{[0,\tau_{K}]}(r)dr - g_{r}\mathbf{1}_{[0,\tau_{K}]}(r)\eta(dr,X_{r}) + Z_{r}\mathbf{1}_{[0,\tau_{K}]}(r) dW_{r},$ and
\begin{align*}
J_{1} &:= \sum_{[t_{i},t_{i+1}]\in \pi} \Bigg[\phi(t_{i+1},Y^{\tau_{K}}_{t_{i+1}}) - \phi(t_{i+1},Y^{\tau_{K}}_{t_{i}}) - \int_{t_{i}}^{t_{i+1}}\partial_{y}\phi(t_{i+1},Y^{\tau_{K}}_{r}) dY^{\tau_{K}}_{r}\\
&\qquad  - \frac{1}{2}\int_{t_{i}}^{t_{i+1}}\partial^{2}_{yy}\phi(t_{i+1},Y^{\tau_{K}}_{r})|Z_{r}|^{2}\mathbf{1}_{[0,\tau_{K}]}(r)dr\Bigg],\\
J_{2} &:= \frac{1}{2}\sum_{[t_{i},t_{i+1}]\in \pi}\left[\int_{t_{i}}^{t_{i+1}}\partial^{2}_{yy}\phi(t_{i+1},Y^{\tau_{K}}_{r})|Z_{r}|^{2}\mathbf{1}_{[0,\tau_{K}]}(r)dr - \int_{t_{i}}^{t_{i+1}}\partial^{2}_{yy}\phi(r,Y^{\tau_{K}}_{r})|Z_{r}|^{2}\mathbf{1}_{[0,\tau_{K}]}(r)dr\right],\\
J_{3} &:= \sum_{[t_{i},t_{i+1}]\in\pi}\left[\int_{t_{i}}^{t_{i+1}}\big(\partial_{y}\phi(t_{i+1},Y^{\tau_{K}}_{r}) - \partial_{y}\phi(r,Y^{\tau_{K}}_{r})\big) dY^{\tau_{K}}_{r}\right].
\end{align*}
Thus, to prove the desired result, it suffices to prove $|J_1|+|J_2| +|J_3|$ converges to 0 in probability as $|\pi|$ tends to 0. For $J_{2}$, due to the uniform continuity of $\partial^{2}_{yy}\phi(r,y)$ on compact sets and the mean value theorem, we have 
$\lim_{|\pi|\downarrow 0}J_{2} = 0,\ a.s.$  The proof of the convergence of $J_1$, $J_3$ will be separated into the following two steps.

\textbf{Step 1.}
In this step, we shall show that $\lim\limits_{|\pi|\downarrow 0}J_{3} = 0$ a.s.  Note that
\begin{equation*}	\begin{aligned}
J_{3} &= \sum_{[t_{i},t_{i+1}]\in \pi} \int_{t_{i}}^{t_{i+1}}\big(\partial_{y}\phi(r,Y^{\tau_{K}}_{r}) - \partial_{y}\phi(t_{i+1},Y^{\tau_{K}}_{r})\big) f_r \mathbf{1}_{[0,\tau_{K}]}(r)dr\\
&\quad + \sum_{[t_{i},t_{i+1}]\in \pi} \int_{t_{i}}^{t_{i+1}}\big(\partial_{y}\phi(t_{i+1},Y^{\tau_{K}}_{r}) - \partial_{y}\phi(r,Y^{\tau_{K}}_{r})\big) Z_r \mathbf{1}_{[0,\tau_{K}]}(r)dW_{r}\\
&\quad +\sum_{[t_{i},t_{i+1}]\in \pi}
\int_{t_{i}}^{t_{i+1}}\big(\partial_{y}\phi(r,Y^{\tau_{K}}_{r}) - \partial_{y}\phi(t_{i+1},Y^{\tau_{K}}_{r})\big) g_r \mathbf{1}_{[0,\tau_{K}]}(r) \eta(dr,X_{r}) =: K_1 + K_2 +K_3.
\end{aligned}
\end{equation*}
The first summation $K_1$ tends to zero a.s. as $|\pi|\to 0$ due to the uniform continuity of $\partial_y \phi(t, y)$ on compact sets. For $K_{2}$, It\^o's isometry yields that
\begin{equation*}		\begin{aligned}	\mathbb{E}\left[|K_{2}|^{2}\right] = \sum_{[t_{i},t_{i+1}]\in\pi} \E\left[\int_{t_{i}}^{t_{i+1}} |\partial_{y}\phi(t_{i+1},Y^{\tau_{K}}_{r})-\partial_{y}\phi(r,Y^{\tau_{K}}_{r})|^{2}|Z_{r}|^{2}\mathbf{1}_{[0,\tau_{K}]}(r) dr\right].
\end{aligned}
\end{equation*}
Then, by the dominated convergence theorem, we have
$\lim_{|\pi|\downarrow 0}\mathbb{E}\left[|K_{2}|^{2}\right] = 0.$			
As for $K_{3}$, denote $p := p_{2}\vee p_{3}$ and  $l := \sup_{[t_{i},t_{i+1}]\in \pi}\left\{i;\ t_{i} < \tau_{K}\right\}$, we have
\begin{align}\label{e:k2}\nonumber
|K_{3}|
&\le \sum_{[t_{i},t_{i+1}]\in \pi} \left\|\int_{\cdot}^{t_{i+1}\land \tau_{K}}(\partial_{y}\phi(t_{i+1}\land \tau_{K},Y^{\tau_{K}}_{r}) - \partial_{y}\phi(r,Y^{\tau_{K}}_{r}))g_{r}\eta(dr,X_{r})\right\|_{\frac{1}{\tau}\text{-var};[t_{i}\land \tau_{K},t_{i+1}\land \tau_{K}]}\\  
&\qquad + \left\|\int_{\cdot}^{\tau_{K}}(\partial_{y}\phi( \tau_{K},Y^{\tau_{K}}_{r}) - \partial_{y}\phi(t_{l+1},Y^{\tau_{K}}_{r}))g_{r}\eta(dr,X_{r})\right\|_{\frac{1}{\tau}\text{-var};[t_{l},\tau_{K}]}.
\end{align}
By \eqref{e:p-var1}, the second term on the right-hand side of \eqref{e:k2} is bounded by  (up to a multiplicative constant)
\begin{align*}						
&\|\eta\|^{(K)}_{\tau,\lambda}\left(1+\|X^{\tau_{K}}_{\cdot}\|^{\lambda}_{p_{1}\text{-var};[t_{l},\tau_K]}\right)|\tau_K-t_l|^\tau \Big(\|(\partial_{y}\phi(\tau_{K},Y^{\tau_{K}}_{\cdot}) - \partial_{y}\phi(t_{l+1},Y^{\tau_{K}}_{\cdot}))g_{\cdot}\|_{\infty;[t_{l},\tau_{K}]}\\
&\qquad  + \|(\partial_{y}\phi(\tau_{K},Y^{\tau_{K}}_{\cdot}) - \partial_{y}\phi(t_{l+1},Y^{\tau_{K}}_{\cdot}))g_{\cdot}\|_{p\text{-var};[t_{l},\tau_{K}]}\Big)\\
&\quad \lesssim  \|\eta\|^{(K)}_{\tau, \lambda} (1+K^{\lambda})|\pi|^{\tau }\Big( K\sup_{r\in[0,T], |y|\le K}|\partial_{y}\phi(r,y)| + \sup_{r\in[0,T]}\| \partial_{y}\phi(r,Y^{\tau_{K}}_{\cdot})g_{\cdot}\|_{p\text{-var};[t_{l},\tau_{K}]}\Big),
\end{align*}
which converges to 0 as $|\pi|\to0$ in view of the uniform boundedness of $\| \partial_{y}\phi(r,Y^{\tau_{K}}_{\cdot})g_{\cdot}\|_{p\text{-var};[t_{l},\tau_{K}]}$. For the first term on the right-hand side of \eqref{e:k2}, by \eqref{e:p-var1} and the fact that 
\begin{align*}		
\big(\partial_{y}\phi(t_{i+1}\land{\tau_{K}},Y^{\tau_{K}}_{r}) - \partial_{y}\phi(r,Y^{\tau_{K}}_{r})\big)\big|_{r=t_{i+1}\land\tau_{K}}=0, 
\end{align*}
it is bounded by (up to a multiplicative constant)
\begin{equation}\label{e:partial phi1}
\|\eta\|^{(K)}_{\tau, \lambda} (1+K^{\lambda}) \sum_{[t_i, t_{i+1}]\in\pi} |t_{i+1} - t_{i}|^{\tau} \big\|\big(\partial_{y}\phi(t_{i+1},Y^{\tau_{K}}_{\cdot}) - \partial_{y}\phi(\cdot,Y^{\tau_{K}}_{\cdot})\big)g_{\cdot}\big\|_{p\text{-var};[t_{i}\land \tau_{K},t_{i+1}\land \tau_{K}]}.
\end{equation}
The $p$-variation above can be bounded as follows,
\begin{align}	\nonumber
\big\|&(\partial_{y}\phi(t_{i+1},Y^{\tau_{K}}_{\cdot}) - \partial_{y}\phi(\cdot,Y^{\tau_{K}}_{\cdot}))g_{\cdot}\big\|_{p\text{-var};[t_{i}\land \tau_{K},t_{i+1}\land \tau_{K}]}\\ \nonumber
&\le 	\big\|(\partial_{y}\phi(t_{i+1},Y^{\tau_{K}}_{\cdot}) - \partial_{y}\phi(\cdot,Y^{\tau_{K}}_{\cdot}))\big\|_{p_{3}\text{-var};[t_{i}\land\tau_{K},t_{i+1}\land\tau_{K}]}\|g_{\cdot}\|_{\infty ;[0,t\land\tau_{K}]}\\ \label{e:partial phi2}
&\quad +\big\|(\partial_{y}\phi(t_{i+1},Y^{\tau_{K}}_{\cdot}) - \partial_{y}\phi(\cdot,Y^{\tau_{K}}_{\cdot}))\big\|_{\infty;[t_{i}\land\tau_{K},t_{i+1}\land\tau_{K}]}\|g_{\cdot}\|_{p_{2}\text{-var} ;[0,t\land\tau_{K}]}\\ 
\nonumber
& \le \big\|(\partial_{y}\phi(t_{i+1},Y^{\tau_{K}}_{\cdot}) - \partial_{y}\phi(\cdot,Y^{\tau_{K}}_{\cdot}))\big\|_{p_{3}\text{-var};[t_{i}\land\tau_{K},t_{i+1}\land\tau_{K}]}\Big(\|g_{\cdot}\|_{\infty ;[0,t\land\tau_{K}]}  + \|g_{\cdot}\|_{p_{2}\text{-var} ;[0,t\land\tau_{K}]} \Big).
\end{align}
We further estimate $	\big\|(\partial_{y}\phi(t_{i+1},Y^{\tau_{K}}_{\cdot}) - \partial_{y}\phi(\cdot,Y^{\tau_{K}}_{\cdot}))\big\|_{p_{3}\text{-var};[t_{i}\land\tau_{K},t_{i+1}\land\tau_{K}]}$ in the above inequality.  For any partition $\tilde \pi$ on $[t_i, t_{i+1}]$, we have
\begin{equation}\label{e:partial phi3}
\begin{aligned}
\bigg\{&\sum_{[s,u]\in \tilde \pi}\Big|\partial_{y}\phi(t_{i+1},Y^{\tau_{K}}_{u}) - \partial_{y}\phi(u,Y^{\tau_{K}}_{u}) - \partial_{y}\phi(t_{i+1},Y^{\tau_{K}}_{s}) + \partial_{y}\phi(s,Y^{\tau_{K}}_{s})\Big|^{p_{3}}\bigg\}^{\frac{1}{p_{3}}}\\
&\le \bigg\{\sum_{[s,u]\in \tilde \pi}\Big|\partial_{y}\phi(t_{i+1},Y^{\tau_{K}}_{u}) - \partial_{y}\phi(t_{i+1},Y^{\tau_{K}}_{s}) - \partial_{y}\phi(u,Y^{\tau_{K}}_{u}) + \partial_{y}\phi(u,Y^{\tau_{K}}_{s})\Big|^{p_{3}}\bigg\}^{\frac{1}{p_{3}}}\\
&\quad+ \bigg\{\sum_{[s,u]\in \tilde \pi}\Big|-\partial_{y}\phi(u,Y^{\tau_{K}}_{s}) + \partial_{y}\phi(s,Y^{\tau_{K}}_{s})\Big|^{p_{3}}\bigg\}^{\frac{1}{p_{3}}}\\
&\le 2\sup_{r\in[0,t],|y|\le K}|\partial^{2}_{yy}\phi(r,y)|\cdot \|Y^{\tau_{K}}_{\cdot}\|_{p_{3}\text{-var};[t_{i},t_{i+1}]} + \|\phi\|^{(K)}_{\gamma,1}\cdot |t_{i+1}-t_{i}|^{\gamma}.
\end{aligned}
\end{equation}
In view of \eqref{e:partial phi1}, \eqref{e:partial phi2}, and \eqref{e:partial phi3}, the first term on the right-hand side of \eqref{e:k2} is bounded by (up to a multiplicative constant that depends on $K$)
\begin{align*}
\left(|\pi|^{\tau} +  \sum_{[t_{i},t_{i+1}]\in\pi}|t_{i+1}-t_{i}|^{\tau}\Big(\|Y^{\tau_{K}}_{\cdot}\|_{p_{3}\text{-var};[t_{i},t_{i+1}]} +|t_{i+1}-t_{i}|^{\gamma}\Big)\right),
\end{align*}
which converges to 0 as $|\pi|\to0$ by Lemma~\ref{lem:control's property}.  
Combining the above analysis on the right-hand side of \eqref{e:k2}, we have
$\lim_{|\pi|\downarrow 0}|K_{3}|  = 0,\  a.s.$ Therefore, we have $\lim_{|\pi|\downarrow 0} J_3 = 0$ a.s.

\textbf{ Step 2.} In this step, we will prove $\lim_{|\pi|\downarrow 0} |J_1|=0$. Taylor's expansion yields that
\begin{equation*}					
\begin{aligned}
J_{1} &= \sum_{[t_{i},t_{i+1}]\in\pi}\Big[ \partial_{y}\phi(t_{i+1},Y^{\tau_{K}}_{t_{i}})\left(Y^{\tau_{K}}_{t_{i+1}}-Y^{\tau_{K}}_{t_{i}}\right) - \int_{t_{i}}^{t_{i+1}}\partial_{y}\phi(t_{i+1},Y^{\tau_{K}}_{r})dY^{\tau_{K}}_{r}\Big]\\
&\qquad + \sum_{[t_{i},t_{i+1}]\in\pi}\Big[ \frac{1}{2} \partial^{2}_{yy}\phi(t_{i+1},\xi_{i})\left(Y^{\tau_{K}}_{t_{i+1}}-Y^{\tau_{K}}_{t_i}\right)^{2}  - \frac{1}{2}\int_{t_{i}}^{t_{i+1}}\partial^{2}_{yy}\phi(t_{i+1},Y^{\tau_{K}}_{r})|Z_{r}|^{2}\mathbf{1}_{[0,\tau_{K}]}(r)dr\Big]\\[1mm]
&=: L_{1} + L_{2},
\end{aligned}
\end{equation*}
where $\xi_{i}$ is a (random) number between $Y_{t_{i}}^{\tau_K}$ and $Y_{t_{i+1}}^{\tau_K}$. 
By a similar argument for $J_3$, we can show that $L_1$ converges to 0 in probability as $|\pi|\to0$. As for $L_{2}$, denote $Y^{\tau_{K}}_{t_{i+1}}-Y^{\tau_{K}}_{t_{i}} =: \alpha_i+\beta_i+\gamma_i$,
where $\alpha_i =  - \int_{t_i}^{t_{i+1}} f_{r} \mathbf{1}_{[0,\tau_{K}]}(r) dr, \beta_i= - \int_{t_i}^{t_{i+1}}g_{r} \mathbf{1}_{[0,\tau_{K}]}(r)\eta(dr,X_{r})$, and $\gamma_i= \int_{t_i}^{t_{i+1}} \mathbf{1}_{[0,\tau_{K}]}(r)Z_{r}dW_{r}.$ Then we have	
\begin{equation}\label{e:estimate-Ito}	
\begin{aligned}
|L_{2}| &\lesssim \sum_{[t_{i},t_{i+1}]\in\pi} \frac{1}{2} |\partial^{2}_{yy}\phi(t_{i+1},\xi_{i})|\left( \alpha_i^{2} +\left|\alpha_i\gamma_i\right| +\beta_i^{2} + \left|\beta_i\gamma_i\right| + |\alpha_i \beta_i|\right)\\
&\quad +\sum_{[t_{i},t_{i+1}]\in\pi}\left| \frac{1}{2} \partial^{2}_{yy}\phi(t_{i+1},\xi_{i}) \gamma_i^{2} - \frac{1}{2} \partial^{2}_{yy}\phi(t_{i+1},Y^{\tau_{K}}_{t_{i}}) \gamma_i^{2}\right|\\						
&\quad + \sum_{[t_{i},t_{i+1}]\in\pi}\left|\frac{1}{2} \partial^{2}_{yy}\phi(t_{i+1},Y^{\tau_{K}}_{t_{i}})  \gamma_i^{2}  - \frac{1}{2}\int_{t_{i}}^{t_{i+1}}\partial^{2}_{yy}\phi(t_{i+1},Y^{\tau_{K}}_{r})|Z_{r}|^{2}\mathbf{1}_{[0,\tau_{K}]}(r)dr\right|.
\end{aligned}
\end{equation}
We claim that $\sum_{[t_{i},t_{i+1}]\in\pi} \left( \alpha_i^{2} +\left|\alpha_i\gamma_i\right| +\beta_i^{2} + \left|\beta_i\gamma_i\right| + |\alpha_i \beta_i|\right)$ converges to 0 in probability as $|\pi|\to0$. Indeed, by Proposition~\ref{prop:bounds}, we have
\begin{align*}
\sum_{i}\left|\beta_i \gamma_i\right|		& \le \sum\limits_{i}\left\|\int_{\cdot}^{t_{i+1}}g_{r}\mathbf{1}_{[0,\tau_{K}]}(r)\eta(dr,X_{r})\right\|_{\frac{1}{\tau}\text{-var};[t_{i},t_{i+1}]} \left\|\int_{\cdot}^{t_{i+1}}Z_{r}\mathbf{1}_{[0,\tau_{K}]}(r)dW_{r}\right\|_{p_{3}\text{-var};[t_{i},t_{i+1}]}\\
&\lesssim_{\tau,\lambda,p_{1},p_{2}}\|\eta\|^{(K)}_{\tau,\lambda}(1+K^{\lambda})\cdot K\cdot \sum_{i}|t_{i+1}-t_{i}|^{\tau}\left\|\int_{\cdot}^{t_{i+1}}Z_{r}\mathbf{1}_{[0,\tau_{K}]}(r)dW_{r}\right\|_{p_{3}\text{-var};[t_{i},t_{i+1}]},
\end{align*}
which converges to 0 a.s. as $|\pi|\to 0$ by Lemma~\ref{lem:control's property} and the fact that $\tau + \frac{1}{p_{3}} > 1$. Similar results hold for the other terms. For the second summation in \eqref{e:estimate-Ito}, we have	
\begin{equation*}
\sum_{i}\frac{1}{2}\left|\partial^{2}_{yy}\phi(t_{i+1},\xi_{i}) - \partial^{2}_{yy}\phi(t_{i+1},Y^{\tau_K}_{t_{i}})\right|\gamma_i^{2}\le \frac{1}{2}\sup_{i}\left|\partial^{2}_{yy}\phi(t_{i+1},\xi_{i})-\partial^{2}_{yy}\phi(t_{i+1},Y^{\tau_K}_{t_{i}})\right| \sum_{i}\gamma_i^{2},
\end{equation*}
where the right-hand side converges to 0 in probability due to the uniform continuity of $\partial^2_{yy}\phi$ on compact sets, 
and the fact that $\sum_{i}\gamma_i^{2}$ converges in probability  to $\int_{0}^{t} |Z_{r}|^{2}\mathbf{1}_{[0,\tau_{K}]}(r) dr$ as $|\pi|\to0$.
Denote by $\frac12H$ the third summation in \eqref{e:estimate-Ito}, we have 
\begin{align*}
H =& \sum_{[t_{i},t_{i+1}]\in\pi}\left(\partial^{2}_{yy}\phi(t_{i+1},Y^{\tau_{K}}_{t_{i}}) \int_{t_{i}}^{t_{i+1}}|Z_{r}|^{2}\mathbf{1}_{[0,\tau_{K}]}(r)dr - \int_{t_{i}}^{t_{i+1}}\partial^{2}_{yy}\phi(t_{i+1},Y^{\tau_{K}}_{r})|Z_{r}|^{2}\mathbf{1}_{[0,\tau_{K}]}(r)dr\right)	\\
&\quad + \sum\limits_{[t_{i},t_{i+1}]\in\pi}\frac{1}{2} \partial^{2}_{yy}\phi(t_{i+1},Y^{\tau_{K}}_{t_{i}}) \left[\left(\int_{t_{i}}^{t_{i+1}}Z_{r}\mathbf{1}_{[0,\tau_{K}]}(r)dW_{r}\right)^{2} - \int_{t_{i}}^{t_{i+1}}|Z_{r}|^{2}\mathbf{1}_{[0,\tau_{K}]}(r)dr\right]\\[1mm]
=: & \ H_{1} + H_{2}.
\end{align*}
Similar to the second summation on the right-hand side of \eqref{e:estimate-Ito}, we have $H_1$ and $H_2$ converging to zero in probability, and thus $\lim_{|\pi|\downarrow 0}|J_1| = 0$ in probability.

\ \ \ \ \

In view of \eqref{decomI}, \eqref{e:I2}, and \eqref{e:I1}, letting $|\pi|\rightarrow 0$, we have
\begin{equation*}
\begin{aligned}
\phi(t,Y_{t\land \tau_{K}}) =   \phi(0,Y_{0}) &+ \int_{0}^{t}  \phi(dr,Y^{\tau_{K}}_{r}) -\int_{0}^{t\land\tau_{K}}\partial_{y}  \phi(r,Y^{\tau_{K}}_{r}) f_{r} dr - \int_{0}^{t\land\tau_{K}} \partial_{y}  \phi(r,Y^{\tau_{K}}_{r})g_{r}\eta(dr,X_{r})\\
&+ \int_{0}^{t\land\tau_{K}}  \partial_{y} \phi(r,Y^{\tau_{K}}_{r})Z_{r} dW_{r} + \frac{1}{2}\int_{0}^{t\land\tau_{K}}\partial^{2}_{yy}  \phi(r,Y^{\tau_{K}}_{r})|Z_{r}|^{2} dr.
\end{aligned}
\end{equation*}
Finally, our desired formula follows if we let $K$ tend to infinity. 
\end{proof}

\end{document}
